\documentclass{beamer}
\usetheme{Berkeley}

\usepackage{
  amsmath,mathtools,amssymb,amsthm,
  centernot,float,diagbox,adjustbox,
  graphicx,caption,subcaption,listings,
  enumerate,xcolor,hyperref,color,soul,
  stmaryrd,amsrefs,tikz,mathrsfs,etoolbox,
  refcount,multicol,MnSymbol,yfonts,framed,
  relsize,colortbl,pgf-umlsd,varwidth,
  environ,xparse,forest}
\usepackage[many]{tcolorbox}
\forestset{qtree/.style={for tree={parent anchor=south, 
           child anchor=north,align=center,inner sep=0pt}}}
\usepackage[framemethod=tikz]{mdframed}
\newcommand{\grayleftbar}[1]{\textcolor[gray]{0.7}{%
     \begin{leftbar}\textcolor{black}{#1}\end{leftbar}}}   
\newcommand{\darkleftbar}[1]{\textcolor[gray]{0.3}{%
     \begin{leftbar}\textcolor{black}{#1}\end{leftbar}}}   
\newcounter{countitems}
\newcounter{nextitemizecount}
\newcommand{\setupcountitems}{%
  \stepcounter{nextitemizecount}%
  \setcounter{countitems}{0}%
  \preto\item{\stepcounter{countitems}}%
}
\makeatletter
\newcommand{\computecountitems}{%
  \edef\@currentlabel{\number\c@countitems}%
  \label{countitems@\number\numexpr\value{nextitemizecount}-1\relax}%
}
\newcommand{\nextitemizecount}{%
  \getrefnumber{countitems@\number\c@nextitemizecount}%
}
\newcommand{\previtemizecount}{%
  \getrefnumber{countitems@\number\numexpr\value{nextitemizecount}-1\relax}%
}

\usepackage[utf8]{inputenc}
\makeatother    
\newenvironment{AutoMultiColItemize}{%
\ifnumcomp{\nextitemizecount}{>}{3}{\begin{multicols}{2}}{}%
\setupcountitems\begin{itemize}}%
{\end{itemize}%
\unskip\computecountitems\ifnumcomp{\previtemizecount}{>}{3}{\end{multicols}}{}}
\usetikzlibrary{
    arrows,shapes,positioning,decorations.markings,
    decorations.pathreplacing,hobby,automata,tikzmark}
\definecolor{mycolor}{rgb}{0.122, 0.435, 0.698}

\newmdenv[innerlinewidth=0.5pt, roundcorner=4pt,linecolor=mycolor,innerleftmargin=6pt,
innerrightmargin=6pt,innertopmargin=6pt,innerbottommargin=6pt]{mybox}

\usepackage[pdf]{graphviz}
\usetikzlibrary{shapes,arrows,calc}
\usepackage[utf8]{inputenc}
\newtheorem{remark}{Remark}

\newcommand{\TODO}[0]{\textcolor{red}{TODO}}
\newcommand{\F}[0]{\mathsf{F}}
\newcommand{\U}[0]{\mathsf{U}}
\newcommand{\G}[0]{\mathsf{G}}
\newcommand{\X}[0]{\mathsf{X}}
\newcommand{\acs}[0]{\textsc{ACL2S}}
\newcommand{\acl}[0]{\textsc{ACL2}}
\newcommand{\pr}[0]{\textsc{Promela}}
\newcommand{\ko}[0]{\textsc{Korg}}
\newcommand{\spn}[0]{\textsc{SPIN}}
\newcommand{\ie}{\emph{i.e.}}
\newcommand{\etal}{\emph{et al.}}

\newcommand{\Pcl}[0]{\textsc{P}}
\newcommand{\NP}[0]{\textsc{NP}}
\newcommand{\SAT}[0]{\texttt{SAT}}
\newcommand{\coNP}[0]{\textsc{coNP}}
\newcommand{\Time}[0]{\text{TIME}}
\newcommand{\Space}[0]{\text{SPACE}}
\newcommand{\threeSAT}[0]{\texttt{3SAT}}

\DeclarePairedDelimiter\abs{\lvert}{\rvert}
\DeclarePairedDelimiter\norm{\lVert}{\rVert}

\usepackage{tikz}
\usetikzlibrary{arrows,chains,matrix,positioning,scopes}

\newlength{\bubblewidth}
\AtBeginDocument{\setlength{\bubblewidth}{.55\textwidth}}
\definecolor{bubblegreen}{RGB}{103,184,104}
\definecolor{bubblegray}{RGB}{241,240,240}
\newcommand{\bubble}[4]{%
  \tcbox[
    colback=#1,
    colframe=#1,
    #2,
  ]{\color{#3}\begin{varwidth}{\bubblewidth}#4\end{varwidth}}%
}

\ExplSyntaxOn
\seq_new:N \l__ooker_bubbles_seq
\tl_new:N \l__ooker_bubbles_first_tl
\tl_new:N \l__ooker_bubbles_last_tl

\NewEnviron{rightbubbles}
 {
  \raggedleft\sffamily
  \seq_set_split:NnV \l__ooker_bubbles_seq { \par } \BODY
  \int_compare:nTF { \seq_count:N \l__ooker_bubbles_seq < 2 }
   {
    \bubble{bubblegreen}{rounded~corners}{white}{\BODY}
   }
   {
    \seq_pop_left:NN \l__ooker_bubbles_seq \l__ooker_bubbles_first_tl
    \seq_pop_right:NN \l__ooker_bubbles_seq \l__ooker_bubbles_last_tl
    \bubble{bubblegreen}{sharp~corners=southeast}{white}{\l__ooker_bubbles_first_tl}\par
    \seq_map_inline:Nn \l__ooker_bubbles_seq
     {
      \bubble{bubblegreen}{sharp~corners=east}{white}{##1}\par
     }
    \bubble{bubblegreen}{sharp~corners=northeast}{white}{\l__ooker_bubbles_last_tl}\par
   }
 }
\NewEnviron{leftbubbles}
 {
  \raggedright\sffamily
  \seq_set_split:NnV \l__ooker_bubbles_seq { \par } \BODY
  \int_compare:nTF { \seq_count:N \l__ooker_bubbles_seq < 2 }
   {
    \bubble{bubblegray}{rounded~corners}{black}{\BODY}
   }
   {
    \seq_pop_left:NN \l__ooker_bubbles_seq \l__ooker_bubbles_first_tl
    \seq_pop_right:NN \l__ooker_bubbles_seq \l__ooker_bubbles_last_tl
    \bubble{bubblegray}{sharp~corners=southwest}{black}{\l__ooker_bubbles_first_tl}\par
    \seq_map_inline:Nn \l__ooker_bubbles_seq
     {
      \bubble{bubblegray}{sharp~corners=west}{black}{##1}\par
     }
    \bubble{bubblegray}{sharp~corners=northwest}{black}{\l__ooker_bubbles_last_tl}\par
   }
 }
\ExplSyntaxOff


\title{Communication Complexity}
\author{Jake Kinsella and Max von Hippel}
\institute{Northeastern University}
\date{\today}



\begin{document}

\frame{\titlepage}

\begin{frame}
\frametitle{Communication Complexity}
\emph{If Alice knows $x$, and Bob knows $y$, how many bits of information must they communicate, in order for both Alice and Bob to know $f(x, y)$?}
\end{frame}

\begin{frame}
\tableofcontents
\end{frame}

% 0  - 5m   - Max introduces the problem
% 5  - 10m  - Jake a couple examples
% 10 - 30m  - Jake does fooling set
% 30 - 50m  - Max does the tiling method
% 50 - 60m  - Max does 2-party discrepancy method
% 60 - 65m  - Jake introduces motivating problem for multi-party version
% 65 - 70m  - Jake defines generalizes multi-party version of problem 
% 70 - 85m  - Max gives multi-party discrepancy without proof
% 85 - 90m  - Jake does non-deterministic version
% 90 - 95m  - Max does randomized

\section{Introduction}

\begin{frame}{Introduction (Max)}
Consider a two-party communication problem,
	in which the participants

\begin{figure}[h]
\centering
\begin{subfigure}{.3\textwidth}
  \centering
  \includegraphics[width=.2\linewidth]{introduction/alice.png}
  \caption{Alice}
  \label{fig:alice}
\end{subfigure}%
and
\begin{subfigure}{.3\textwidth}
  \centering
  \includegraphics[width=.2\linewidth]{introduction/bob.png}
  \caption{Bob}
  \label{fig:bob}
\end{subfigure}
\label{fig:participants}
\end{figure}

\emph{participate} to compute a function:
\[f : \underbrace{\mathbb{B}^n}_{\begin{aligned}\text{Alice's}\\\text{input}\end{aligned}} \times 
      \underbrace{\mathbb{B}^n}_{\begin{aligned}\text{Bob's}\\\text{input}\end{aligned}} \to 
      \underbrace{\mathbb{B}}_{\begin{aligned}\text{global}\\\text{output}\end{aligned}}\]
\end{frame}

\begin{frame}{Introduction (Max)}
The players can come up with a \emph{protocol} 
\(\Pi = (p_1, ..., p_t)\),
namely, 
	for some natural $t \in \mathbb{N}$,
	a sequence of $t$-many functions
\(p_i : \mathbb{B}^* \to \mathbb{B}^*\)
such that the communication between the players looks like this ...
\end{frame}

\begin{frame}{Introduction (Max)}
{\color{blue}{\footnotesize Alice is given input $x$.}}
\pause
\begin{leftbubbles}
{\footnotesize Hello Bob.  I can't reveal $x$, but $p_1(x)$ is \texttt{p1}.}
\end{leftbubbles}
\pause
\rightline{{\color{olive}{\footnotesize Bob is given input $y$.}}}
\pause
\begin{rightbubbles}
{\footnotesize Thanks Alice.  I can't reveal $y$, but $p_2(y, \texttt{p1})$ is \texttt{p2}.}
\end{rightbubbles}
\pause
\center{\footnotesize{... yada yada yada ...}}
\pause
\begin{leftbubbles}
{\footnotesize Pleasure doing business with you Bob.  My final clue for you is that $p_{n-1}(x, \texttt{p1}, ..., \texttt{pn-2})$ is \texttt{pn-1}.}
\end{leftbubbles}
\pause
\begin{rightbubbles}
{\footnotesize Rad.  Then $p_n(y, \texttt{p1}, ..., \texttt{pn-1})$ is \texttt{pn}.  TTFN!}
\end{rightbubbles}
\end{frame}

\begin{frame}{Introduction (Max)}
\begin{itemize}
\item The functions $p_i$ can be \emph{anything} so long as they are well-defined.  E.g., could solve the Halting Problem.
\item After the final message, \emph{both parties} must know $f(x, y)$.
\end{itemize}
\pause
\begin{definition}[Communication Complexity]
Suppose $\Pi$ is a protocol for $f$ in which at most $t$
bits are communicated between Alice and Bob.
Then the \emph{communication complexity} of $\Pi$,
	denoted $C(\Pi)$, is $t$.
\end{definition}
\pause
\begin{definition}[$C(f)$]
The communication complexity of $f$, denoted $C(f)$,
	is the minimum communication complexity
	achieved by any protocol for $f$.
\end{definition}
\end{frame}

\subsection{Examples}

\begin{frame}{Parity (Jake)}
\begin{example}[Are the number of \texttt{1}s in $xy$ even (0), or odd (1)?]
\(f : \mathbb{B}^n \times \mathbb{B}^n \to \mathbb{B}\)
is precisely
\((x, y) \mapsto \bigoplus xy\).
\end{example}
\pause
Example protocol $\Pi$:
\begin{leftbubbles}
{\footnotesize $\texttt{P1}=\text{parity}(x)$.}
\end{leftbubbles}
\pause
\begin{rightbubbles}
{\footnotesize $\texttt{P2}=\text{parity}(y) \bigoplus \texttt{P1}$}
\end{rightbubbles}
\pause
Now both Alice and Bob know $f(x, y) = \texttt{P2}$.
$C(f) \leq 2$ because $C(\Pi) = 2$ and $\Pi$ implements $f$.
But $C(f) \geq 2$ because $f$ depends on $x$ and $y$.
Hence \(C(f) = 2\).
\end{frame}

\begin{frame}{Halting (Jake)}
\begin{example}[$A_{TM}$]
\(H : \mathbb{B}^n \times \mathbb{B}^n \to \mathbb{B}\)
is precisely
\(\langle M, 1^{n} \rangle \mapsto 1\)
if $M$ halts on $1^{n}$ else 0.
\end{example}
\pause
Example protocol $\Pi$:\\
\begin{rightbubbles}
{\footnotesize $\texttt{P1}=1 \text{ if } y = 1^{n} \text{ else } 0$.}
\end{rightbubbles}
\pause
\begin{leftbubbles}
\texttt{P2}$=(M$ does/doesn't halt on $1^{n})$.
\end{leftbubbles}
\pause
Both players have unlimited computation power.
We are only interested in communication complexity.
\end{frame}

\section{Methods}

\subsection{2-Party Problem}

\begin{frame}{Methods (Max)}
If we find a protocol $\Pi$,
	then we know $C(f)$ is at most $C(\Pi)$.
\pause

What if we don't know any protocol $\Pi$?
\pause
\begin{itemize}
	\item Could we upper-bound $C(f)$ without knowing $\Pi$?
\end{itemize}

\pause
What if the only protocols we find seem really lousy?
\pause
\begin{itemize}
	\item Could we lower-bound $C(f)$ without finding a better protocol?
\end{itemize}

\pause
\centering{\textbf{TL;DR: yup.}}
\end{frame}

\subsubsection{Fooling Set Method}

\begin{frame}{Fooling Set Method (Jake)}
We begin with a motivating observation.
\pause
\begin{lemma}[Communication Equality is Image Equality]
If Alice and Bob exchange the same sequence of messages when Alice gets $\textbf{x}$ and Bob gets $\textbf{y}$ as they do when Alice gets $\textbf{x}'$ and Bob gets $\textbf{y}'$, then $f(\textbf{x}, \textbf{y}) = f(\textbf{x}', \textbf{y}')$.
\end{lemma}
\pause
\begin{proof}
$\Pi$ is deterministic and $f$ is a function.
\end{proof}
\pause
\underline{Idea:} an efficient protocol will efficiently group together inputs that go to the same output.

\end{frame}

\begin{frame}{Fooling Set Method (Jake)}
\begin{definition}[Fooling Set]
If $f : \mathbb{B}^n \times \mathbb{B}^n \to \mathbb{B}$ is a function,
	a \emph{fooling set} for $f$ is a set
	\(S \subseteq \mathbb{B}^n \times \mathbb{B}^n\)
	such that for some choice $b \in \mathbb{B}$
\(f(S) = \{ b \}\)
but, for all distinct $(x, y), (x', y') \in S$,
either $f(x, y') \neq b \text{ or } f(x', y) \neq b$
\end{definition}
Basically, a fooling set is a group of inputs that go to the same output,
but which is \emph{brittle} to argument-swapping.
In some sense these \emph{brittle} sets lower-bound the difficulty in grouping like inputs.
\pause
\begin{lemma}[Fooling Set Method]
If $f$ has a size-$M$ fooling set, then $C(f) \geq \text{log}_2(M)$.
\end{lemma}
\end{frame}

\begin{frame}{Fooling Set Method (Jake)}
\begin{example}[Set-Disjointness]
$\textsc{Disj} : \mathbb{B}^n \times \mathbb{B}^n \to \mathbb{B}$ is the function
that maps $(A, B)$ to 1 if $A \cap B = \emptyset$ else 0.
\end{example}

How many fooling sets does \textsc{Disj} have?
\pause
Notice $A \cap B = \emptyset$ if $B = \overline{A}$.
\pause
There are $2^n$ possible values $A$.
\pause
None of these distinct $(A, \overline{A}), (A', \overline{A'})$ satisfy
 $A \cap \overline{A} = A \cap \overline{A'}$ or
$A \cap \overline{A} = A' \cap \overline{A}$
else they wouldn't be distinct.
\pause
So we get a $2^n$-size fooling set.
\pause
\[\therefore \, C(\textsc{Disj}) \geq \text{log}_2(2^n) = n\]
\end{frame}

\begin{frame}{Fooling Set Method (Jake)}
\textbf{NTS:} If $f$ has a size-$M$ fooling set then $C(f) \geq \text{log}_2(M)$.
\begin{proof}
For a contradiction suppose a protocol $\Pi$ exists for $f$ s.t. $C(\Pi) < \text{log}_2(M)$.
\pause
Then $\Pi$ yields at most $2^{C(\Pi)} < 2^{\text{log}_2(M)} = M$
distinct communication patterns.
\pause
However, there are $M$ input pairs $(x, y)$ that each have distinct communication patterns by definition.
\pause
Since $2^{C(\Pi)} < M$ there must be some $(x, y), (x', y')$ on which $\Pi$ yields the same communication pattern.

\medskip

\pause
Then $(x, y')$ must yield the same communication pattern as $(x, y)$
as Bob cannot possibly tell the difference.  The argument is symmetric for $(x', y)$ and Alice.  One of the two must yield a contradiction and we are done.
\end{proof}
\end{frame}

\subsubsection{Tiling Method}

\begin{frame}{Tiling Method (Max)}
With the \emph{fooling set} method, we lower-bounded $C(f)$.
Now we'll introduce a new method that both lower- and upper-bounds $C(f)$.
\pause
\begin{definition}[$M(f)$]
The \emph{matrix of }$f$,
	denoted $M(f)$,
		is the $2^n \times 2^n$ matrix whose $(x, y)$th entry
		is the value $f(x, y)$.
\end{definition}
\pause
\begin{example}[$M(\lor)$]
\adjustbox{max width=.25\textwidth}{
	\begin{tabular}{llllll}
	   & {\cellcolor{blue!20}{00}} & 
	     {\cellcolor{blue!20}{01}} & 
	     {\cellcolor{blue!20}{10}} & 
	     {\cellcolor{blue!20}{11}} \\
	{\cellcolor{green!20}{00}} & 00 & 01 & 10 & 11 \\
	{\cellcolor{green!20}{01}} & 01 & 01 & 11 & 11 \\
	{\cellcolor{green!20}{10}} & 10 & 11 & 10 & 11 \\
	{\cellcolor{green!20}{11}} & 11 & 11 & 11 & 11
	\end{tabular}
}
\adjustbox{max width=.73\textwidth}{
\begin{minipage}{\linewidth}
\begin{itemize}
	\item The green cells are Alice's possible inputs $x$.
	\item The blue cells are Bob's possible inputs $y$.
	\item The uncolored cells are the matrix $M(f)$.
\end{itemize}
\end{minipage}
}
\end{example}
\end{frame}

\begin{frame}{Tiling Method (Max)}
\begin{definition}[Combinatorial Rectangle]
A \emph{combinatorial rectangle} in $M(f)$ is any submatrix of $M$.
We say a rectangle $A \times B$ in $M(f)$ is \emph{monochromatic}
if for all $x, x'$ in $A$ and $y, y'$ in $B$, $M_{x,y} = M_{x',y'}$.
\end{definition}
\pause
\medskip

\underline{Idea:} Each event in a protocol $\Pi$ splits $M(f)$ into two or more
combinatorial rectangles of still-possible values for $f(x, y)$.
\pause
\medskip

\underline{Intuition:} Much like splitting a circuit $C$ into ``$C$ where the first bit is 0'' and ``$C$ where the first bit is 1''.
\pause
\medskip

Let's see an example ...
\end{frame}

\begin{frame}{Tiling Method (Max)}
\begin{example}[$\Pi=\textsc{LeastSignificantBit},f=<$]
\begin{itemize}
\item $f : \mathbb{B}^3 \times \mathbb{B}^3 \to \mathbb{B}$ is the function that maps $(x, y)$ to 1 if $x < y$ else 0.  
\item $\Pi = \textsc{LeastSignificantBit}$ is the na{\"i}ve protocol where Alice and Bob read off their bits from right to left.
\end{itemize}
\adjustbox{max width=.8\textwidth}{
\centering
\begin{forest}
[{\begin{tabular}{|l|lllll|}
\hline
    & \texttt{000}
    & \texttt{001}
    & \texttt{010}
    & \texttt{011} 
    & \texttt{100}\\\hline
\texttt{000} & \cellcolor{red!25}0 & 
                     \cellcolor{green!25}1 & 
                     \cellcolor{green!25}1 & 
                     \cellcolor{green!25}1 & 
                     \cellcolor{green!25}1 \\
\texttt{001} & \cellcolor{red!25}0 & 
                     \cellcolor{red!25}0 & 
                     \cellcolor{green!25}1 & 
                     \cellcolor{green!25}1 & 
                     \cellcolor{green!25}1 \\
\texttt{010} & \cellcolor{red!25}0 & 
                     \cellcolor{red!25}0 & 
                     \cellcolor{red!25}0 & 
                     \cellcolor{green!25}1 & 
                     \cellcolor{green!25}1 \\
\texttt{011} & \cellcolor{red!25}0 & 
                     \cellcolor{red!25}0 & 
                     \cellcolor{red!25}0 & 
                     \cellcolor{red!25}0 & 
                     \cellcolor{green!25}1 \\
\texttt{100} & \cellcolor{red!25}0 & 
                     \cellcolor{red!25}0 & 
                     \cellcolor{red!25}0 & 
                     \cellcolor{red!25}0 & 
                     \cellcolor{red!25}0\\\hline
\end{tabular}},
	[{\begin{tabular}{|l|lllll|}
\hline
    & \texttt{000}
    & \texttt{001}
    & \texttt{010} 
    & \texttt{011}
    & \texttt{100}\\\hline
\texttt{000}  & \cellcolor{red!25}0 & 
                     \cellcolor{green!25}1 & 
                     \cellcolor{green!25}1 & 
                     \cellcolor{green!25}1 & 
                     \cellcolor{green!25}1 \\
\texttt{010}  & \cellcolor{red!25}0 & 
                     \cellcolor{red!25}0 & 
                     \cellcolor{red!25}0 & 
                     \cellcolor{green!25}1 & 
                     \cellcolor{green!25}1 \\
\texttt{100}  & \cellcolor{red!25}0 & 
                     \cellcolor{red!25}0 & 
                     \cellcolor{red!25}0 & 
                     \cellcolor{red!25}0 & 
                     \cellcolor{red!25}0\\\hline
\end{tabular}}, edge label={node[midway,left]{Alice: ``$x=\_\_\texttt{0}$''}}]
	[{\begin{tabular}{|l|lllll|}
\hline
    & \texttt{000} 
    & \texttt{001}
    & \texttt{010}
    & \texttt{011}
    & \texttt{100}\\\hline
\texttt{001}  & \cellcolor{red!25}0 & 
                     \cellcolor{red!25}0 & 
                     \cellcolor{green!25}1 & 
                     \cellcolor{green!25}1 & 
                     \cellcolor{green!25}1 \\
\texttt{011}  & \cellcolor{red!25}0 & 
                     \cellcolor{red!25}0 & 
                     \cellcolor{red!25}0 & 
                     \cellcolor{red!25}0 & 
                     \cellcolor{green!25}1 \\\hline
\end{tabular}}, edge label={node[midway,right]{Alice: ``$x=\_\_\texttt{1}$''}}]]
\end{forest}}
\end{example}
\end{frame}

\begin{frame}{Tiling Method (Max)}
Now we get to the punchline.
\pause
\begin{definition}[Monochromatic Tiling]
A \emph{monochromatic tiling} of $M(f)$ is a partition of $M(f)$ into disjoint
monochromatic rectangles.
\end{definition}
\pause
\begin{center}
\textbf{It's thinking time.}
\end{center}
\pause
\begin{itemize}
\item Then the leaves of the tree induced by $\Pi$ and rooted at $M(f)$ clearly form a monochromatic tiling of $M(f)$.
\pause
\item The number of leaves in a binary tree can be used to upper-bound its depth.
\pause
\item The depth of the binary tree induced by $\Pi$ is exactly $C(\Pi)$.
\end{itemize}
\end{frame}

\begin{frame}{Tiling Method (Max)}
\begin{center}
\textbf{It's thinking time.}
\end{center}
\begin{itemize}
\item Then the leaves of the tree induced by $\Pi$ and rooted at $M(f)$ clearly form a monochromatic tiling of $M(f)$.
\item The number of leaves in a binary tree can be used to upper-bound its depth.
\item The depth of the binary tree induced by $\Pi$ is exactly $C(\Pi)$.
\end{itemize}
\pause
\begin{theorem}[The Punchline]
Let $\chi(f)$ denote the minimum number of rectangles in any monochromatic tiling of $M(f)$.
\[\text{log}_2 \chi(f) \leq
	C(f) \leq
		16 \big( \text{log}_2 \chi(f) \big)^2
\]
\end{theorem}
\end{frame}

\begin{frame}{Tiling Method (Max)}
\textbf{NTS:} $\text{log}_2 \chi(f) \leq C(f)$.
\begin{proof}
Assume $C(f)$. 
\pause
Then $\exists$ a protocol $\Pi$ in which $\leq C(f)$ bits are communicated
between the 2 participants. 
\pause 
For simplicity suppose each bit is communicated individually. 
\pause 
Then $\Pi$ induces a tree whose max depth is $C(f)$, whose leaves form a monochromatic partition of $M(f)$.
\pause  
Every m.c. partition / $M(f)$ requires $\geq \chi(f)$ rectangles, so the tree induced by $\Pi$ has $\geq \chi(f)$ leaves.
\pause  
But it's a binary tree so its depth is at least $\text{log}_2 \chi(f)$.
\end{proof}
\end{frame}

% This proof is *enormous*, and appears to contain
% a small arithmetic error somewhere near the end 
% (a missing 2 coefficient).  So, we probably should not
% present it, but there's no harm in including it in the
% supplemental materials and sharing it online.
\begin{frame}{Tiling Method (Max)}
{\small \textbf{NTS:} $C(f) \leq 16 \big( \text{log}_2 \chi(f) \big)^2$.~\cite{aho1983notions}}
\begin{proof}
\only<1>{
\begin{itemize}
	\item Let $M_1, ..., M_{\chi(f)}$ be a monochromatic
	partitioning of $M(f)$ known ahead of time to both
	Alice (on the ``left'') 
	and Bob (on the ``right'').
	Each rectangle $M_i$ can alternatively be written $X_i \times Y_i$.
	\item Let $G_L, G_R$ be graphs whose nodes are $\{ 1, ..., \chi(f) \}$.
	There is an edge $i \to j$ in $G_L$ ($G_R$ resp.)
		if $M_i$ and $M_j$ have a row (column resp.) in common.
	\item Let $\text{deg}_L(u)$ (resp. $\text{deg}_R(u)$)
		denote the degree of the node $u$ in the graph $G_L$
		(resp. $G_R$.)
	\item Let $x$ be Alice's input and $y$ Bob's input.
\end{itemize}
}
\only<2>{
\begin{itemize}
	\item We'll describe the protocol in ``rounds''.
	During the rounds, Alice keeps track of $Y$ (a set containing $y$)
	and Bob keeps track of $X$ (a set containing $x$),
	both of which are initially $\mathbb{B}^n$.
	\item Both sides know the graphs $G_L, G_R$ and the rectangles $M_i$
	ahead of time.
\end{itemize}
}
\only<3>{
Each stage proceeds as follows.
\begin{enumerate}
	\item Alice looks for a rectangle $M_i = X_i \times Y_i$ s.t. $x \in X_i$ and $\text{deg}_L(i) \leq 3\chi(f)/4$.
	\begin{enumerate}
		\item If she finds some such rectangle then she sends $i$ to Bob.
		\begin{enumerate}
			\item Bob replies to indicate if $y \in M_i$.  
			\item If so then the protocol ends because $f(x, y)$ is the color of $M_i$.
			\item Otherwise $X := X \cap X_i$, and each rectangle $M_{\alpha} = X_{\alpha} \cap Y_{\alpha}$ is replaced with $(X_i \cap X_{\alpha}) \times Y_{\alpha}$.
		\end{enumerate}
		\item Otherwise she replies that she found no such rectangle.
		In this case Bob does what Alice just attempted, symmetrically, with a small caveat ...
	\end{enumerate}
\end{enumerate}
}
\only<4>{s
\begin{itemize}
\item If neither Alice nor Bob could find any $M_i$ with low-enough degree,
	then they both know that every node $i$ in $G := G_L \cap G_R$
	for which $(x, y) \in M_i$
	has degree $\geq (3\chi(f)/4)^2$ $= 9\chi(f)/16$ $> \chi(f)/2$.
\item Let $i, j$ both have degree $\geq \chi(f)/2$ in $G$.
	Then some node $z$ is adjascent to $i$ and $j$ in $G$, 
	by the Pigeonhole Principle.
	Hence $M_i \cap M_z \neq \emptyset$ and $M_j \cap M_z \neq \emptyset$.
	But the rectangles are monochromatic, hence, $M_i$ and $M_j$ are the same color.  So Alice needs to find an $M_i$ containing $x$ and some $y \in Y$ whose degree in $G$ is at least $\chi(f)/2$; and Bob's procedure is symmetric.
\end{itemize}
}
\only<5>{
\begin{itemize}
	\item In the worst case for each stage, the first participant sends ``nothing found'' (1 bit), the second participant sends some $i$ ($\leq \text{log}_2(\chi(f))$ bits), and the first participant replies with some $j$ ($\leq \text{log}_2(\chi(f))$ bits).  So in the worst case each round requires $\leq 1 + 2\text{log}_2(\chi(f))$ bits.
	\item The protocol ends after at most $n$ rounds where $(3\chi(f)/4)^n \approx 1$, i.e., after $\text{log}_{(4/3)}(\chi(f))$ rounds.
	\item So total communication complexity is $\leq \text{log}_{(4/3)}(\chi(f)) * (1 + 2\text{log}_2(\chi(f)))$.
\end{itemize}
}
\only<6>{
	For $\chi(f) \geq 2$:
	\[\begin{aligned}
	C(f) & \leq \text{log}_{(4/3)}(\chi(f)) * (1 + 2\text{log}_2(\chi(f))) \\
		 & = \frac{\text{log}_2(\chi(f))}{\text{log}_2(4/3)} * (1 + 2\text{log}_2(\chi(f))) \\
		 & < 2.5 * \text{log}_2(\chi(f)) * (1 + 2\text{log}_2(\chi(f))) \\
		 & \leq 2.5 * \text{log}_2(\chi(f)) * 3\text{log}_2(\chi(f))) \\
		 & = 7.5 \text{log}_2^2(\chi(f)) \\
		 & \leq 16 \text{log}_2^2(\chi(f))
	\end{aligned}\]
	{\tiny I'm almost certainly missing a factor of 2, which would explain the choice of 16. -Max.}
}
\alt<6>{\qedhere}{\phantom\qedhere}
\end{proof}

\end{frame}

\subsubsection{Discrepency Method}

\begin{frame}{2-Party Discrepency Method (Max)}
When we partition $M(f)$ into some number of rectangles,
	the sizes of the rectangles must add up to the size of $M(f)$.
\pause
\medskip

Hence, if $\chi(f) \leq K$ for some integer $K$,
	then $M(f)$ must have a m.c. rectangle containing at least $2^n * 2^n / K$ entries.
\pause
\begin{proof}
Suppose $\chi(f) \leq K$ for some integer $K$.
\pause
If $\chi(f) = K$ then $\exists$ a partioning of $M(f)$ into $K$ m.c. rects, in which case at least 1 must have size $\geq \abs{M(f)} / K$, i.e., $2^n * 2^n / K$.  
\pause
On the other hand if $\chi(f) < K$ then $\chi(f) = K'$ for some $K' < K$ and then $M(f)$ can be partitioned into $K'$ monochromatic rectangles, at least 1 of which has size $\geq \abs{M(f)} / K'$, which is strictly larger than $\abs{M(f)} / K$.
\pause
Either way the conjecture holds.
\end{proof}
\end{frame}

\begin{frame}{2-Party Discrepency Method (Max)}
Suppose that $M(f)$ contains a monochromatic rectangle $A \times B$ having at least $2^n * 2^n / K$ entries.
\pause
Since $A \times B$ is monochromatic, this implies that:
\[\sum_{a \in A, b \in B} (-1)^{M_{a,b}} = \begin{cases}
-1 * \abs{A \times B} & \text{ if it's colored } 1 \\
+1 * \abs{A \times B} & \text{ if it's colored } 0\end{cases}\]
\pause
So if we wrap an absolute value above our sum, we get:
\[\abs{\sum_{a \in A, b \in B} (-1)^{M_{a,b}}} = \text{ the size of the rectangle } A \times B\]
\pause
But we already assumed that $A \times B$ has at least $2^n * 2^n / K$ entries, hence:
\[\abs{\sum_{a \in A, b \in B} (-1)^{M_{a,b}}} = \geq 2^n * 2^n / K\]
\end{frame}

\begin{frame}{2-Party Discrepency Method (Max)}
Let's divide both size by $2^n * 2^n$, for fun and profit.
\[\frac{1}{2^n * 2^n}\abs{\sum_{a \in A, b \in B} (-1)^{M_{a,b}}} \geq 1 / K\]
\pause
Mathematicians like to name things.
\begin{definition}[Discrepency]
The \textit{discrepency} of a rectangle $A \times B$ of $M(f)$ is exactly the following.
\[\text{Disc}(A \times B) = \frac{1}{2^n * 2^n}\abs{\sum_{a \in A, b \in B} (-1)^{M_{a,b}}}\]
The \emph{discrepency of $M(f)$} is the max disc among its rectangles.
\end{definition}
\end{frame}

\begin{frame}{2-Party Discrepency Method (Max)}
Now that we've named this thing, let's re-write our inequality.
\[\text{Disc}(A \times B) \geq 1 / K\]
\pause
Taking inverses:
\[\frac{1}{\text{Disc}(A \times B)} \leq K\]
\pause
Certainly $\chi(f) \leq \chi(f)$, so supplanting $\chi(f)$ for $K$ in the statement, we get:
\begin{lemma}[2-Party Discrepency Method]
\[\frac{1}{\text{Disc}(A \times B)} \leq \chi(f)\]
\end{lemma}
\end{frame}

\subsection{Multi-Party Problem}

\begin{frame}{Multi-Party Problem (Jake)}
\TODO
\end{frame}

\subsubsection{Multi-Party Discrepency Method}

\begin{frame}{Multi-Party Discrepency Method (Max)}
\TODO
\end{frame}

\section{Other Variants}

\subsection{Non-Deterministic}

\begin{frame}{Non-Deterministic (Jake)}
\TODO
\end{frame}

\subsection{Randomized}

\begin{frame}{Randomized (Max)}
\TODO
\end{frame}

\section{References}

\begin{frame}[t, allowframebreaks]{References}
\bibliographystyle{amsalpha}
\bibliography{pres}
\end{frame}

\end{document}